\documentclass[a4paper,10.5pt]{article}

\usepackage[margin=0.5in,top=0.35in,bottom=0.35in]{geometry}
\usepackage{titlesec}
\usepackage{enumitem}
\usepackage[hidelinks]{hyperref}
\usepackage{fancyhdr}
\usepackage{setspace}
\usepackage{tabularx}

\pagestyle{empty}
\setstretch{0.96}

%---------- Formatting ----------
\titleformat{\section}{\large\bfseries\scshape}{}{0em}{}[\titlerule]
\titlespacing*{\section}{0pt}{4pt}{2pt}
\setlist[itemize]{noitemsep, topsep=0pt, parsep=0pt, left=12pt}

%---------- Document ----------
\begin{document}

%---------- Header ----------
\begin{center}
    {\LARGE \textbf{R Sarang}}\\[3pt]
    \href{mailto:sarangraghavan14@gmail.com}{sarangraghavan14@gmail.com} $\mid$ +91-94980 88240 $\mid$ Chennai, India \\
    \href{https://github.com/rsarang14}{github.com/rsarang14} $\mid$ \href{https://www.linkedin.com/in/sarang-raghavan-a04a2b308}{linkedin.com/in/R SARANG}
\end{center}

%---------- Education ----------
\section*{Education}

\noindent
\begin{tabularx}{\linewidth}{@{}X r@{}}
\textbf{Indian Institute of Information Technology, Design and Manufacturing (IIITDM), Kancheepuram} & \textit{May 2028} \\
\end{tabularx}\\[1pt]
Dual Degree: B.Tech in Electronics \& Communication and M.Tech (Honors) in \textbf{VLSI Design} \hfill \textit{CGPA: 9.01}\\
\textit{Relevant Coursework:} \textbf{Digital IC Design}, \textbf{CMOS VLSI System Design}, \textbf{VLSI Testing \& Testable Design (DFT)}, \textbf{Microprocessors}


\noindent
\begin{tabularx}{\linewidth}{@{}X r@{}}
\textbf{Indian Institute of Technology (IIT), Madras (Online)} & \textit{May 2028} \\
\end{tabularx}\\[1pt]
BS in \textbf{Data Science and Applications} \hfill \textit{CGPA: 8.48}\\
\textit{Relevant Coursework:} Linear Algebra, Optimization, Probabilistic Models, Machine Learning

%---------- Skills ----------
\section*{Technical Skills}
\textbf{Hardware \& VLSI:} \textbf{Verilog}, \textbf{SystemVerilog}, \textbf{RTL Design}, \textbf{Digital ASIC Design}, \textbf{Design for Test (DFT)}, Fault Detection, \textbf{CPU Design}, \textbf{RISC-V Architecture}, CMOS Logic.\\
\textbf{EDA Tools \& Simulation:} QuestaSim, Xilinx Vivado, Intel Quartus Prime, Cadence (Genus, Conformal, Innovus, Virtuoso).\\
\textbf{Software \& Scripting:} C, C++, Python, \textbf{Tcl}, Bash/Shell Scripting, Linux (Ubuntu), Git, \textbf{Object-Oriented Programming}.\\
\textbf{Embedded Systems:} STM32, ESP32, Raspberry Pi, NVIDIA Jetson, Sensor Integration (ADC/DAC, I2C, SPI).

%---------- Professional Development & Certifications ----------
\section*{Certifications \& Awards}
\begin{itemize}[leftmargin=1.5em, itemsep=1pt]
    \item \textbf{ISWDP Certification} -- Advanced training in semiconductor process technology, device design, and microelectronics R\&D tools \hfill \textit{2026}
    \item \textbf{Highest CGPA Award (Ranked 1st in VLSI)} -- IIITDM Kancheepuram \hfill \textit{2023-24, 2024-25}
    \item \textbf{Hardware \& Autonomy Awards:} 1st Place at AtomQuest (2025) $\mid$ Judges Innovation Award, Caterpillar Autonomy Challenge, IITM (2026)
\end{itemize}

%---------- Projects ----------
\section*{Core Hardware \& VLSI Projects}

\noindent
\begin{tabularx}{\linewidth}{@{}X r@{}}
\textbf{1 TOPS (1 Tape Out Per Student) -- VLSI Society of India (VSI) \& C-DAC} & \textit{Nov 2025 -- Present} \\
\end{tabularx}
\begin{itemize}[leftmargin=1.5em, itemsep=1pt]
    \item Selected for a highly competitive national cohort to design and tape-out a \textbf{22nm TSMC RISC-V System-on-Chip (SoC)} targeted for assistive audio signal processing, integrating a \textbf{RISC-V (Ibex) CPU core} over a custom \textbf{AHB} interconnect.
    \item \textbf{Digital ASIC / RTL design} of a dedicated \textbf{256-point FFT/IFFT hardware accelerator} in \textbf{Verilog/SystemVerilog} for real-time frequency-domain noise suppression, using a resource-shared (folded) butterfly datapath over 16-bit fixed-point complex arithmetic to minimize silicon area.
    \item Architected a \textbf{zero-copy, ping-pong RAM datapath} enabling the FFT, IFFT, and CPU filtration stages to operate in-place on shared SRAM, cutting originally planned DMA transfers by \textbf{$\sim$80\%}, backed by a round-robin \textbf{AHB arbiter/decoder} with sticky-grant arbitration across CPU and DMA masters.
    \item Verified full-chip functionality via an \textbf{800M-cycle (100 MHz) SystemVerilog testbench} exercising four consecutive real-time audio frames end-to-end, with in-sim cycle-count instrumentation to benchmark CPU vs. DMA vs. accelerator performance.
\end{itemize}


\noindent
\begin{tabularx}{\linewidth}{@{}X r@{}}
\textbf{Custom RISC-V CPU on FPGA -- eYRC 2025 (MazeSolver)} & \textit{Sep 2025 -- Mar 2026} \\
\end{tabularx}
\begin{itemize}[leftmargin=1.5em, itemsep=1pt]
    \item Designed and verified a custom \textbf{32-bit RISC-V CPU core} from scratch using \textbf{Verilog} (\textbf{CPU Design}, \textbf{RTL Design}).
    \item Synthesized and implemented the RTL on an \textbf{Intel Quartus FPGA}, achieving a maximum clock frequency of \textbf{150 MHz}.
    \item Engineered hardware-level control units to handle real-time I/O sensor processing, validating closed-loop PID motor actuation on an isolated single-motor testbench with zero OS overhead.
\end{itemize}

%---------- Experience ----------
\section*{Experience}

\noindent
\begin{tabularx}{\linewidth}{@{}X r@{}}
\textbf{Controls Engineering Intern -- Kineshia Robotics} & \textit{May 2026 -- Aug 2026} \\
\end{tabularx}
\begin{itemize}[leftmargin=1.5em, itemsep=1pt]
    \item Engineered prehensile walking controllers and complex null-space projection algorithms for the Graspman open-chain robotic manipulator using \textbf{C++/Python} (\textbf{Object-Oriented Programming}), increasing operational workspace and grasping success rates by 20\%.
    \item Executed system-level hardware planning, developing 50+ comprehensive CAD models and finalizing the Bill of Materials (BoM), accelerating physical prototyping cycles by 3 weeks.
    \item Optimized Prioritized Inverse Kinematics (IK) control loops, reducing computational latency by 40\% (to under 1ms) to ensure real-time stability and high-frequency (1 kHz) trajectory execution.
\end{itemize}

%---------- Leadership ----------
\section*{Leadership \& Collaboration}

\noindent
\begin{tabularx}{\linewidth}{@{}X r@{}}
\textbf{Student Technical Manager \& Co-Lead -- MaRS (Mars Rover Students) Club} & \textit{2025-2026} \\
\end{tabularx}
\begin{itemize}[leftmargin=1.5em, itemsep=1pt]
    \item Led technical hardware integration and mentored a multidisciplinary engineering team of \textbf{40} students.
    \item Presented system architecture and embedded design workflows to \textbf{NASA JPL scientists} at IEEE WAMS 2025.
\end{itemize}

\end{document}
